\chapter{GIỚI THIỆU ĐỀ TÀI}

Chương này trình bày bối cảnh hình thành đề tài, mục tiêu cần đạt được, phạm vi triển khai và định hướng giải pháp của UniTrack. Nội dung chương làm rõ bài toán quản lý tiến độ dự án sinh viên trước khi chuyển sang phân tích yêu cầu ở Chương 2.

\section{Đặt vấn đề}

Trong các học phần đồ án, nghiên cứu khoa học sinh viên hoặc đồ án tốt nghiệp, sản phẩm cuối cùng chỉ phản ánh một phần kết quả đánh giá. Giảng viên còn phải theo dõi dự án đang hoạt động, nhiệm vụ đã giao, sinh viên phụ trách, bài nộp đang chờ phản hồi, nhóm có nguy cơ chậm tiến độ và minh chứng được dùng khi đánh giá. Khi những thông tin này nằm rải rác trong email, nhóm chat, bảng tính và thư mục lưu file, quá trình giám sát phụ thuộc nhiều vào ghi nhớ thủ công.

Quá trình thực hiện dự án thường có một vòng lặp lặp lại. Giảng viên giao việc trong buổi họp hoặc qua tin nhắn; sinh viên gửi liên kết, ảnh chụp, báo cáo hoặc sản phẩm qua nhiều kênh; phản hồi lại xuất hiện trong một luồng trao đổi khác. Sau một thời gian, cả hai phía phải rà soát lại để xác định phiên bản mới nhất, bài đã được duyệt và bài cần sửa. Hạn chế chủ yếu nằm ở việc nhiệm vụ, bài nộp, minh chứng và quyết định đánh giá chưa được đặt trong cùng một ngữ cảnh dữ liệu.

UniTrack được xây dựng từ nhu cầu trên. Hệ thống lựa chọn hướng tiếp cận lấy dự án làm trung tâm: mỗi dự án có giảng viên phụ trách, thành viên sinh viên, mốc kế hoạch, nhiệm vụ chính thức, bài nộp, tài nguyên tham khảo, file minh chứng và lịch sử duyệt bài. Phân hệ chính của sản phẩm được đặt tên là \textit{Workspace}, nhấn mạnh phạm vi sử dụng là không gian làm việc của dự án thay vì mạng xã hội, kho học liệu độc lập hoặc hệ thống chấm điểm số riêng.

\tableintro{tab:student-project-pain-points}{tổng hợp các khó khăn nghiệp vụ khi giảng viên và sinh viên theo dõi dự án bằng nhiều công cụ rời rạc}
\begin{table}[H]
\caption{Các khó khăn nghiệp vụ khi giám sát dự án sinh viên bằng công cụ rời rạc}
\label{tab:student-project-pain-points}
\centering
\begin{tblr}{
  width=\textwidth,
  colspec={X[0.95,l]X[1.45,l]X[1.45,l]},
  rows={valign=t},
  row{1}={font=\bfseries,bg=black!7},
  hlines={0.35pt,black!24},
}
Tình huống & Cách làm rời rạc thường gặp & Vấn đề UniTrack cần xử lý \\
Giao nhiệm vụ & Giảng viên nhắn trong nhóm chat hoặc ghi trong bảng tính & Nhiệm vụ thiếu trạng thái rõ ràng, khó gắn với mốc kế hoạch và người được giao \\
Nộp tiến độ & Sinh viên gửi link hoặc file qua nhiều kênh & Khó biết bài nộp nào đang chờ duyệt, bài nào đã được trả lại để sửa \\
Theo dõi minh chứng & Tệp nằm trong Drive, máy cá nhân hoặc tin nhắn & Minh chứng không gắn chặt với bài nộp và có thể bị sửa sau khi duyệt \\
Quản lý thành viên & Danh sách nhóm được ghi thủ công & Sinh viên không còn active hoặc không thuộc dự án vẫn có thể xuất hiện trong dữ liệu phụ \\
Tổng hợp tiến độ & Giảng viên rà từng dự án bằng tay & Dễ bỏ sót dự án có bài chờ duyệt, nhiệm vụ quá hạn hoặc nhóm thiếu tiến độ \\
\end{tblr}
\end{table}

\section{Mục tiêu đề tài}

Mục tiêu của đồ án là xây dựng một nền tảng Web hỗ trợ giảng viên giám sát tiến độ dự án sinh viên theo một luồng thống nhất: \textit{Giảng viên -- Dự án -- Nhiệm vụ -- Bài nộp -- Duyệt bài}. Phạm vi hệ thống tập trung vào những thao tác quan trọng nhất để quản lý dự án và phản hồi tiến độ, chưa mở rộng sang toàn bộ hoạt động đào tạo.

Về chức năng, hệ thống cần cho phép quản trị viên quản lý tài khoản; giảng viên tổ chức Workspace, tạo dự án, quản lý nhóm, lập mốc kế hoạch và giao nhiệm vụ; sinh viên xem việc được giao và nộp báo cáo tiến độ kèm minh chứng; giảng viên duyệt bài nộp và lưu lại phản hồi. Về dữ liệu, các đối tượng quan trọng như dự án, nhiệm vụ, bài nộp, tài nguyên và minh chứng phải được đặt trong cùng một ngữ cảnh để người dùng biết trạng thái mới nhất của công việc. Về an toàn, hệ thống phải kiểm tra quyền theo vai trò, quan hệ với dự án và trạng thái vòng đời của dự án thay vì chỉ dựa vào các nút hiển thị trên giao diện. Về chất lượng, đồ án cần có kiểm thử cho các luồng rủi ro cao như đăng nhập, phân quyền, nộp bài, duyệt bài và bảo toàn minh chứng.

\section{Phạm vi đề tài}

UniTrack hiện tập trung vào quản lý dự án học tập được giảng viên giám sát. Một số tên bảng lịch sử hoặc ý tưởng mở rộng vẫn còn trong migration, nhưng phần sản phẩm đang hoạt động được giữ gọn để tránh biến hệ thống thành nhiều module độc lập không cần thiết.

\tableintro{tab:implementation-scope}{tóm tắt phạm vi triển khai thực tế của UniTrack và phân biệt phần đã hoàn thành với phần còn để phát triển tiếp}
\begin{table}[H]
\caption{Phạm vi triển khai thực tế của UniTrack}
\label{tab:implementation-scope}
\centering
\begin{tblr}{
  width=\textwidth,
  colspec={X[1.0,l]X[2.1,l]X[0.7,c]},
  rows={valign=t},
  row{1}={font=\bfseries,bg=black!7},
  hlines={0.35pt,black!24},
}
Nhóm chức năng & Phạm vi trong phiên bản hiện tại & Trạng thái \\
Xác thực/session & Đăng nhập, đăng xuất, xem người dùng hiện tại, session cookie, giới hạn đăng nhập và thu hồi session khi tài khoản thay đổi & \okbadge \\
Quản lý tài khoản & Tạo tài khoản, tìm kiếm, đổi vai trò/trạng thái, đặt mật khẩu và hướng dẫn chuyển trạng thái khi tài khoản còn công việc mở & \okbadge \\
Workspace/thư mục & Thư mục lớp nhẹ, danh sách dự án, chi tiết thư mục, thêm dự án có sẵn và khớp chủ thư mục/giảng viên phụ trách & \okbadge \\
Dự án/nhóm & Tạo/sửa/xem dự án, trạng thái vòng đời, tổng hợp tiến độ, nhóm, thêm sinh viên active và vai trò leader/member & \okbadge \\
\stackcell{Nhiệm vụ\\bài nộp} & Mốc kế hoạch, nhiệm vụ, người được giao, bài nộp, trang duyệt bài và trạng thái nhiệm vụ suy ra & \okbadge \\
Tài nguyên/minh chứng & Đường dẫn tài nguyên, tải lên/tải xuống/xóa file minh chứng, lưu trữ riêng tư và khóa dữ liệu sau duyệt & \okbadge \\
Nhật ký hoạt động & Ghi log cho quản lý tài khoản và thao tác thành viên dự án; chưa có giao diện audit log & \partialbadge \\
Phân tích quản trị & Chưa có màn hình quản trị toàn bộ dự án hoặc phân tích nâng cao & \futurebadge \\
\end{tblr}
\end{table}

Các chức năng không thuộc phạm vi chính gồm: đăng ký công khai, invite link mở, chat thời gian thực, mạng xã hội, hệ thống khóa học độc lập, chấm điểm số chi tiết, rubric phức tạp, thanh toán, thông báo qua email và phân tích dữ liệu nâng cao. Việc loại bỏ những chức năng này giúp hệ thống tập trung vào luồng giám sát dự án và giữ mã nguồn theo từng lát tính năng rõ ràng.

\section{Định hướng giải pháp}

Hệ thống gồm bốn thành phần chính. Giao diện React hiển thị màn hình, quản lý dữ liệu tạm trên trình duyệt và gọi API. Server Go cung cấp REST API dưới tiền tố \route{/api/v1}, kiểm tra session, quyền truy cập, dữ liệu đầu vào và trạng thái vòng đời. PostgreSQL lưu dữ liệu nghiệp vụ cùng các ràng buộc toàn vẹn. Lớp lưu trữ dùng thư mục cục bộ khi phát triển và Cloudflare R2 private bucket khi triển khai production.

\figurelead{fig:solution-overview.png}{khái quát cách giao diện, API, cơ sở dữ liệu và kho minh chứng phối hợp để tạo thành luồng xử lý thống nhất cho hệ thống}
\figSolutionOverview

Một yêu cầu ghi dữ liệu không mặc nhiên hợp lệ chỉ vì giao diện đang hiển thị nút thao tác. Giao diện có trách nhiệm ẩn hoặc vô hiệu hóa thao tác không hợp lệ; server vẫn là lớp quyết định cuối cùng. Các xử lý quan trọng khóa dòng dự án trong transaction để kiểm tra lại quyền quản lý, trạng thái vòng đời và trạng thái đối tượng mới nhất trước khi ghi dữ liệu. Với những ràng buộc có thể bị phá bởi thao tác đồng thời hoặc ghi trực tiếp vào cơ sở dữ liệu, PostgreSQL tiếp tục bảo vệ bằng khóa ngoại, chỉ mục duy nhất và trigger.

\begin{artifactnote}{Nguyên tắc thiết kế xuyên suốt}
UniTrack kiểm tra đồng thời ba yếu tố: \textit{người dùng là ai}, \textit{người dùng liên quan thế nào với dự án}, và \textit{dự án đang ở trạng thái nào}. Với dữ liệu phụ như nhiệm vụ, bài nộp, đường dẫn tài nguyên hoặc file minh chứng, hệ thống còn xác nhận đối tượng đó thật sự thuộc cùng dự án. Đây là điểm khác biệt quan trọng giữa một ứng dụng chỉ tạo/sửa/xóa dữ liệu đơn giản và một hệ thống quản lý tiến độ có nhiều trạng thái nghiệp vụ.
\end{artifactnote}

\section{Bố cục báo cáo}

Bố cục báo cáo gồm sáu chương. Chương 1 đặt vấn đề từ thực tế giám sát dự án sinh viên, xác định mục tiêu, phạm vi triển khai và nguyên tắc giải pháp của UniTrack. Phần này giúp người đọc nắm rõ lý do hệ thống chọn hướng project-first trước khi đi vào đặc tả chi tiết.

Chương 2 phân tích bài toán nghiệp vụ qua hiện trạng sử dụng công cụ rời rạc, nhóm tác nhân, use case trọng tâm, quy trình nghiệp vụ và các yêu cầu chức năng/phi chức năng. Các nội dung trong chương này là cơ sở để đối chiếu thiết kế và kết quả triển khai ở những chương sau.

Chương 3 giới thiệu các nền tảng lý thuyết và công nghệ được dùng trong đồ án, gồm kiến trúc ứng dụng Web, REST API, session cookie, transaction, ràng buộc dữ liệu và bộ công nghệ React, Go, PostgreSQL. Trọng tâm của chương là giải thích vì sao những công nghệ đó phù hợp với bài toán giám sát dự án có phân quyền và trạng thái nghiệp vụ.

Chương 4 mô tả thiết kế và triển khai hệ thống: kiến trúc tổng quan, request boundary, route/API, gói mã nguồn, mô hình dữ liệu, giao diện, công cụ, kết quả chức năng, kiểm thử và triển khai tham khảo. Chương này cũng cung cấp các bảng, sơ đồ và ảnh giao diện để nối yêu cầu nghiệp vụ với mã nguồn thực tế.

Chương 5 đi sâu vào các giải pháp có giá trị kỹ thuật nổi bật, như session an toàn hơn token lưu trong browser, kiểm soát truy cập theo dự án, khóa vòng đời trong transaction, ràng buộc toàn vẹn dữ liệu ở PostgreSQL và bảo toàn minh chứng sau duyệt. Các phần này làm rõ phần đóng góp riêng của đồ án so với một ứng dụng CRUD đơn giản.

Chương 6 tổng hợp kết quả đạt được, so sánh UniTrack với cách quản lý bằng công cụ rời rạc, nêu các hạn chế còn tồn tại và đề xuất hướng phát triển tiếp theo. Nội dung cuối báo cáo giúp xác định phần nào đã sẵn sàng sử dụng và phần nào cần tiếp tục hoàn thiện trong các phiên bản sau.

Từ phạm vi đã xác định, báo cáo chuyển sang phân tích tác nhân, use case, quy trình nghiệp vụ và yêu cầu hệ thống ở Chương 2.
